Les isotopes d'un radioélément diffèrent par le nombre de nucléons, leurs
propriétés et leur stabilité. Ils jouent un rôle important dans la vie courante
que ce soit en datation, en médecine ou en industrie.

L'or $\ce{_{79}^{198}Au}$ est un radioisotope de l'or, il se désintègre en
mercure $\ce{_{Z}^{A}Hg}$ par émission $\beta^{-}$. Cette radioactivité est
notamment utilisée dans le traitement contre certains cancers ou pour traiter
d'autres maladies.

Pour évaluer l'efficacité thérapeutique de l'or 198, une équipe de recherche a
porté son attention sur un groupe de souris souffrantes d'un cancer de la
prostate sévère en leurs injectant une dose d'or $\ce{_{79}^{198}Au}$.

\begin{donnees}
    \begin{table}[H]
        \centering
        \setlength{\arrayrulewidth}{0.8pt}
        \renewcommand{\arraystretch}{1.5}
        \begin{tabularx}{\textwidth}{|p{4cm}|C|C|C|C|C|}
            \hline
            Noyau ou particule & $\ce{_{79}^{198}Au}$ & $\ce{_{Z}^{A}Hg}$ &
            Proton & Neutron & Électron \\
            \hline
            Masse en (\unit{\u}) & $197,968225$ & $197,966769$ & $1,007276$ &
            $1,008665$ & $\num{5.49e-4}$ \\
            \hline
            \multicolumn{6}{|l|}{\raisebox{2mm}[9mm]{
                Énergie de liaison par nucléon~:
                $\frac{E_{\ell}}{A}\left(\ce{_{Z}^{A}Hg}\right)=
                \qty{8.4}{\MeV\per\mathrm{nucléon}}$}} \\
            \hline
            \multicolumn{2}{|c|}{$\qty{1}{\u}=\qty{931.5}{\U}$} &
            \multicolumn{4}{c|}{$t_{1/2}\left(\ce{_{79}^{198}Au}\right)=
            \qty{2.7}{jours}$} \\
            \hline
        \end{tabularx}
    \end{table}
\end{donnees}

\begin{enumerate}
    \item Définir un noyau radioactif.
    \item Donner la composition du noyau $\ce{_{79}^{198}Au}$.
    \item Écrire l'équation de désintégration de l'or $\ce{_{79}^{198}Au}$ en
          précisant les valeurs de $A$ et $Z$.
    \item Vérifier que la valeur de l'énergie de liaison
          $E_{\ell}(\ce{_{79}^{198}Au})$ est égale à $\qty{1525.53}{\MeV}$.
    \item Justifier que le noyau $\ce{_{Z}^{A}Hg}$ est plus stable que la noyau
          $\ce{_{79}^{198}Au}$.
    \item Calculer, en unité $\unit{\MeV}$, l'énergie libérée
          $E_{\text{libérée}}=|\Delta E|$ par la réaction de désintégration d'un
          noyau d'or 198.
    \item On injecte à une souris à $t_{0}=0$ une dose radioactive d'or 198
          d'activité $a_{0}=\qty{5e6}{\Bq}$.

          Après 28 jours de l'injection, le volume de la tumeur devient plus
          faible.

          Déterminer l'activité a de la dose après 28 jours de l'injection.
\end{enumerate}

